\setcounter{section}{3}

\Needspace{5\baselineskip}
\hypertarget{ux57faux4e8eux53efux9a8cux8bc1ux5956ux52b1ux7684ux5f3aux5316ux5b66ux4e60}{%
\section{基于可验证奖励的强化学习}\label{ux57faux4e8eux53efux9a8cux8bc1ux5956ux52b1ux7684ux5f3aux5316ux5b66ux4e60}}

\Needspace{5\baselineskip}
\hypertarget{ux4e00ux6838ux5fc3ux601dux60f3-6}{%
\subsection{一、核心思想}\label{ux4e00ux6838ux5fc3ux601dux60f3-6}}

基于可验证奖励的强化学习（Reinforcement Learning with Varifiable Reward，RLVR）是强化微调的重要范式，使大模型的推理能力实现了飞跃。对于经过预训练的模型，以数学题的答案、编程题的输出结果等的对错或可自动验证的得分作为奖励，可以激发出模型的逻辑推理能力，让模型自行探索出获得正确答案的策略。

这里原则上不引入LLM评分，以免出现Reward Hacking。

\Needspace{5\baselineskip}
\hypertarget{ux4e8cux63d0ux53d6ux7b54ux6848ux7684ux65b9ux5f0f}{%
\subsection{二、提取答案的方式}\label{ux4e8cux63d0ux53d6ux7b54ux6848ux7684ux65b9ux5f0f}}

模型的输出一般为长思维链，包括思考过程（放在 \texttt{\textless{}think\textgreater{}...\textless{}/think\textgreater{}} 标签中）和最终结果（放在 \texttt{\textless{}answer\textgreater{}...\textless{}/answer\textgreater{}} 标签中）。其中，正确答案默认为最终结果里的最后一个LateX字符串。如果在RL训练初期模型没有遵守格式（例如没有输出 \texttt{\textless{}answer\textgreater{}} 标签），系统会直接判定为错误（Reward=0），或者给予一个极小的惩罚。这会迫使模型为了拿到奖励，必须先学会``遵守格式''，然后再学会``做对题目''。

对于数学题，当答案有多种形式时，可用数学公式符号化验证的工具进行判定。对于编程题，可先运行代码，再根据结果给出奖励。对于一些半开放问题如``写出一种满足条件的化学分子式''，答案可能有无数种，可用Checker Script（专用检查脚本）。

\Needspace{5\baselineskip}
\hypertarget{ux4e09ux53efux8bfbux6027ux63d0ux5347}{%
\subsection{三、可读性提升}\label{ux4e09ux53efux8bfbux6027ux63d0ux5347}}

\begin{enumerate}
\def\labelenumi{\arabic{enumi}.}
\tightlist
\item
  硬规则量化：可读性首先需要格式合规性。
\end{enumerate}

（1）结构完整性：是否包含 \texttt{\textless{}think\textgreater{}} 和 \texttt{\textless{}answer\textgreater{}} 标签？是否闭合？标签内是否有内容？

（2）重复度：用N-gram 重复率计算，纯RL模型在卡住时很容易陷入死循环（如不断重复``所以\ldots 所以\ldots 所以\ldots{}''）。系统检测生成的文本中n-gram的重复频率，如果超过阈值，给予巨大的负奖励。

（3）语言一致性：用语言检测算法（如fastText）的置信度，如果Prompt是中文，而Chain of Thought中间突然变成了英文或乱码（DeepSeek-R1-Zero 早期经常出现中英夹杂），这被视为可读性差。系统会计算目标语言的token占比，占比越低，惩罚越重。

（4）长度：用Token数量统计，防止模型为了凑字数或者过度推导而生成冗长的废话。通常会有一个软性的长度范围，超出部分会带来微小的负奖励。

\Needspace{4\baselineskip}
\begin{enumerate}
\def\labelenumi{\arabic{enumi}.}
\setcounter{enumi}{1}
\tightlist
\item
  模型量化：奖励模型（Reward Model, RM）
\end{enumerate}

对于逻辑是否通顺、排版是否美观、语气是否像人这种高阶可读性，规则写不出来，于是用另一个模型来打分。这其实就是RLHF（人类反馈强化学习）的思路。

在纯推理模型（如o1或R1）的早期训练阶段，这个比重很低，甚至没有，因为过分追求``好读''可能会损害模型的智商（模型会为了讨好人类而简化复杂的推理步骤）。但在后期这比较重要。

\Needspace{4\baselineskip}
\begin{enumerate}
\def\labelenumi{\arabic{enumi}.}
\setcounter{enumi}{2}
\tightlist
\item
  最关键的手段：SFT``冷启动''与其后的KL散度约束
\end{enumerate}

这可能违背直觉，但实际上，大部分RLVR模型的高可读性，不是在RL阶段练出来的，而是从SFT（监督微调）阶段继承来的。在开始RL之前，模型先通过少量高质量的、人类（或强模型）编写的``完美过程''进行微调。这些数据格式完美、逻辑清晰、语言优雅。而RL训练时，系统会计算当前模型（Policy Model）和SFT模型（Reference Model）之间的KL散度。

如果模型为了解题开始胡言乱语（可读性下降），KL 散度会飙升，导致总奖励下降。这实际上是将可读性隐式地量化为了与人类写作习惯的相似度。

\Needspace{4\baselineskip}
\begin{enumerate}
\def\labelenumi{\arabic{enumi}.}
\setcounter{enumi}{3}
\tightlist
\item
  现实案例：DeepSeek-R1-Zero的教训
\end{enumerate}

DeepSeek 的论文中提供了一个非常精彩的反面教材，证明了如果没有专门的可读性约束，RL 会把可读性彻底毁掉。DeepSeek-R1-Zero没有经过SFT冷启动，没有加可读性奖励，纯RL后模型确实变聪明了，做题正确率飙升，但输出变得极度难以阅读。

\FloatBarrier
\Needspace{5\baselineskip}
\hypertarget{ux53c2ux8003ux6587ux732e-82}{%
\subsection{参考文献}\label{ux53c2ux8003ux6587ux732e-82}}

\vspace{0.25em}
\begin{itemize}
\tightlist
\item
  DeepSeek-AI. (2025). \href{https://arxiv.org/abs/2501.12948}{DeepSeek-R1: Incentivizing Reasoning Capability in LLMs via Reinforcement Learning}. arXiv:2501.12948.
\item
  温睦宁、林江浩、张伟楠、俞勇（2025）。\href{https://haa.boyuai.com/}{《动手学大模型智能体》}。人民邮电出版社。ISBN 978-7-115-68638-1。
\end{itemize}

\clearpage
